\answer{ex:19:1}{(a)~시냅스 하나는 $6.7\mV$를 준다($R=66.7$~M$\Omega$). 따라서 문턱값에 아예 도달하려면 최소 $4$개가 필요하다(세 개로는 점근적으로만 도달한다); $10\ms$ 안에는 $8$개; $5\ms$ 안에는 $14$개. (b)~$0.1\ms\ll\tau_m$, 누설에 의한 보정은 $\sim0.25\%$.}
\answer{ex:19:2}{(a)~$1.75\cdot10^{10}$(믹서의 4분의 1이 무엇에나 반응한다), $4.4\cdot10^9$, $0.06$($k=40$에서는 거의 반응하지 않는다). (b)~$10^3$배: $2\cdot10^5$ 대 $200$.}
\answer{ex:19:3}{(a)~이항 계수의 대칭성에 의해 $C(2n,n)=2^{2n-1}$. (b)~$0.93$과 $0.09$; 전이 폭은 $P$에 대해 $\sim\sqrt{2n}$, 상대 폭은 $\sim1/\sqrt n$.}
\answer{ex:19:4}{가지돌기 하나에서는 입력이 몇 개든 $V_{\text{세포체}}<\kappa E_E<\theta$; $g_0=0.5g_L$이면 가지돌기마다 $n\ge3$이 필요하다.}
\answer{ex:19:5}{$I\ge C\sigma_V/\sigma_t=15$~nA, 즉 동기 시냅스 $150$개로 비현실적이다. 작은 뉴런은 $1$~nA면 충분하고, $4$~nA 시냅스에서 지터는 $5$~\textmu s.}
\answer{ex:19:6}{가까운 끝의 최댓값: $0.03\,\tau_m$와 $0.97$($L=0.2$); $0.32\,\tau_m$와 $0.66$($L=1$); $0.79\,\tau_m$와 $0.33$($L=2$). 전체 커패시턴스에 의한 추정~\eqref{eq:dendriteload}은 $L\ll1$에서만 맞다.}
\answer{ex:19:7}{열린 문제. $m$이 고정이면 응답 시간은 기억하는 $P$에 선형으로 늘어난다. 비율 $m=\varphi n$이 고정이면 $n$에 의존하지 않게 되며, 큰 뉴런이 느린 것은 크기 자체가 아니라 많은 입력을 합산하기 때문이다.}
